\documentclass{custom_article}

\usepackage[
  backend=biber,
  style=alphabetic
]{biblatex}
\addbibresource{refs.bib}

\title{Intuitionistic Propositional Logic}
\author{Alexander Read}
\date{\today}

\begin{document}

\maketitle

\section*{Introduction}

We consider the positive implicational fragment of the intuitionistic propositional calculus, a language $\mathcal{L}^{\to}$. The alphabet consists of $\textsf{PV} \equiv \{\ p_{0}, \ p_{1}, \ p_{2}, \dots \}$, an infinite set of propositional variables, and a binary connective `${\to}$'. The well-formed formulae (`wff') are defined by the following grammar, where ${\phi},{\psi}$ are meta-variables ranging over formulae:\footnote{Do we define parentheses as part of the vocabulary? Studd does, but Leader notes that they're not strictly part of the language, just used to disambiguate the ASTs. See also the \emph{Unique Readability Theorem} in the B1 §2 notes. We can also use these notes and the Problem Set 1 to define Polish notation and relate it to the prefix notation.}
%
\[
    {\phi},{\psi} ::= \textsf{PV}\ \vert\ ({\phi} \to {\psi})
\]

\noindent I've defined $\textsf{PV}$ in this way because it closely mirrors my implementation in Haskell. However, in this post I'll use `$p$', `$q$', and `$r$' as variables to match the usual notation.

Anyway, following \L ukasiewicz, Meredith, Prior, and others (e.g., see
\href{https://link.springer.com/chapter/10.1007/978-94-010-9614-0_17}{here}),
a lot of the papers I've been reading on $\mathcal{L}^{\to}$ use prefix, rather
than infix, notation to write formulae. For example, where `$C$' is the
prefix-style implication operator, you often see formulae written using the
style on the left, not that on the right:
%
\begin{align}
    CCpCqrCCpqCpr\ {\equiv}\ &(p \to (q \to r)) \to ((p \to q) \to (p \to r)) \tag{S} \\
    CpCqp\ {\equiv}\         &(p \to (q \to p)). \tag{K}
\end{align}

\noindent I hope you agree that the prefix notation is pretty awkward to read.\footnote{The examples correspond to the types of the combinators $\mathsf{\mathbf{S}}$ and $\mathsf{\mathbf{K}}$ in combinatory logic.}

So, I wanted a program to convert from prefix to infix notation. This is done
by parsing valid prefix strings into ASTs, and then printing those trees in
infix format (i.e., by in-order traversal). Rather than use an existing library
like \href{https://hackage.haskell.org/package/parsec}{Parsec}, this was a nice
opportunity to implement a simple parser by hand, following the tutorials from
Hutton and Meijer
(\href{https://www.cs.nott.ac.uk/~pszgmh/monparsing.pdf}{1996},
\href{https://www.cambridge.org/core/journals/journal-of-functional-programming/article/monadic-parsing-in-haskell/E557DFCCE00E0D4B6ED02F3FB0466093}{1998}).

\end{document}
